%==============================================================================
\section{Thuật toán Naive Bayes}
\label{sec:naive-bayes}
%==============================================================================

\begin{dinhnghia}[title={Naive Bayes là gì?}]
Thuật toán Naive Bayes là thuật toán phân loại dựa trên \textbf{định lý Bayes}.
Thuật toán giả định các feature \textbf{độc lập} với nhau --- vì vậy gọi là ``naive'' (ngây thơ).
Tính xác suất mẫu thuộc một lớp dựa trên xác suất của từng feature.
\end{dinhnghia}

\begin{vidu}[title={Ví dụ điện thoại thông minh}]
Điện thoại có thể được coi là ``thông minh'' nếu có màn hình cảm ứng, internet, camera tốt, v.v.
Dù các feature phụ thuộc nhau trong thực tế, mỗi feature vẫn \textbf{đóng góp độc lập} vào xác suất điện thoại thông minh.
\end{vidu}

\begin{dinhnghia}[title={Xác suất hậu nghiệm trong phân loại Bayes}]
Trong phân loại Bayes, quan tâm chính là xác suất hậu nghiệm: xác suất nhãn $L$ khi đã quan sát feature, $P(L \mid \text{features})$.
Theo định lý Bayes:
\[
P(L \mid \text{features}) =
\frac{P(L)\, P(\text{features} \mid L)}{P(\text{features})}
\]
\end{dinhnghia}

\begin{ynghia}[title={Ý nghĩa các thành phần}]
\begin{itemize}
  \item $P(L \mid \text{features})$: xác suất hậu nghiệm của lớp;
  \item $P(L)$: xác suất tiên nghiệm (prior) của lớp;
  \item $P(\text{features} \mid L)$: likelihood --- xác suất predictor khi biết lớp;
  \item $P(\text{features})$: xác suất tiên nghiệm của predictor.
\end{itemize}
\end{ynghia}

\begin{phuongphap}[title={Quy trình Naive Bayes}]
Dùng định lý Bayes tính xác suất mẫu thuộc từng lớp.
Với mỗi feature, tính $P(\text{feature} \mid L)$ rồi \textbf{nhân} các likelihood (giả định độc lập).
Nhân likelihood với prior $P(L)$ để có posterior.
Lặp cho mọi lớp; chọn lớp có xác suất cao nhất.
\end{phuongphap}

\subsection{Các loại thuật toán Naive Bayes}

\begin{ynghia}[title={Nhiều biến thể}]
Có nhiều loại Naive Bayes; phần này trình bày ba loại phổ biến.
\end{ynghia}

\begin{dinhnghia}[title={Gaussian Naive Bayes}]
Gaussian Naive Bayes là bộ phân loại Naive Bayes đơn giản nhất: giả định dữ liệu mỗi nhãn được rút từ phân phối Gaussian.
Dùng khi feature là biến \textbf{liên tục} tuân theo phân phối chuẩn.
\end{dinhnghia}

\begin{dinhnghia}[title={Multinomial Naive Bayes}]
Giả định feature được rút từ phân phối \textbf{Multinomial}.
Phù hợp feature dạng \textbf{đếm rời rạc}; thường dùng phân loại văn bản (tần suất từ trong tài liệu).
\end{dinhnghia}

\begin{dinhnghia}[title={Bernoulli Naive Bayes}]
Giả định feature \textbf{nhị phân} (0 và 1).
Phân loại văn bản với mô hình ``bag of words'' có thể dùng Bernoulli Naive Bayes.
\end{dinhnghia}

\subsection{Triển khai Naive Bayes bằng Python}

\begin{dongco}[title={Gaussian Naive Bayes}]
Tùy dataset chọn mô hình phù hợp; ở đây triển khai \textbf{Gaussian Naive Bayes} bằng scikit-learn.
\end{dongco}

\begin{vidu}[title={Import}]
\begin{lstlisting}
import numpy as np
import matplotlib.pyplot as plt
import seaborn as sns
sns.set()
\end{lstlisting}
\end{vidu}

\begin{vidu}[title={Tạo blob phân phối Gaussian}]
\begin{lstlisting}
from sklearn.datasets import make_blobs
X, y = make_blobs(300, 2, centers=2, random_state=2,
                  cluster_std=1.5)
plt.scatter(X[:, 0], X[:, 1], c=y, s=50, cmap='summer')
plt.show()
\end{lstlisting}

\tsimg{04_naive_bayes/img01.jpg}
\end{vidu}

\begin{vidu}[title={Huấn luyện GaussianNB}]
\begin{lstlisting}
from sklearn.naive_bayes import GaussianNB
model_GNB = GaussianNB()
model_GNB.fit(X, y)
\end{lstlisting}
\end{vidu}

\begin{vidu}[title={Dự đoán trên dữ liệu mới}]
\begin{lstlisting}
rng = np.random.RandomState(0)
Xnew = [-6, -14] + [14, 18] * rng.rand(2000, 2)
ynew = model_GNB.predict(Xnew)

plt.scatter(X[:, 0], X[:, 1], c=y, s=50, cmap='summer')
lim = plt.axis()
plt.scatter(Xnew[:, 0], Xnew[:, 1], c=ynew, s=20,
            cmap='summer', alpha=0.1)
plt.axis(lim)
plt.show()
\end{lstlisting}

\tsimg{04_naive_bayes/img02.jpg}
\end{vidu}

\begin{vidu}[title={Xác suất hậu nghiệm}]
\begin{lstlisting}
yprob = model_GNB.predict_proba(Xnew)
yprob[-10:].round(3)
\end{lstlisting}

\textbf{Output:}
\begin{verbatim}
array([[0.998, 0.002],
       [1.   , 0.   ],
       [0.987, 0.013],
       [1.   , 0.   ],
       [1.   , 0.   ],
       [1.   , 0.   ],
       [1.   , 0.   ],
       [1.   , 0.   ],
       [0.   , 1.   ],
       [0.986, 0.014]])
\end{verbatim}
\end{vidu}

\subsection{Ưu và nhược điểm phân loại Naive Bayes}

\begin{tinhchat}[title={Ưu điểm (Pros)}]
\begin{itemize}
  \item Dễ triển khai và \textbf{nhanh}.
  \item Hội tụ nhanh hơn mô hình phân biệt như logistic regression.
  \item Cần \textbf{ít dữ liệu huấn luyện}.
  \item Mở rộng tuyến tính theo số predictor và số điểm dữ liệu.
  \item Dự đoán xác suất; xử lý dữ liệu liên tục và rời rạc.
  \item Dùng cho phân loại \textbf{nhị phân} và \textbf{đa lớp}.
\end{itemize}
\end{tinhchat}

\begin{luuy}[title={Nhược điểm (Cons)}]
\begin{itemize}
  \item Giả định \textbf{độc lập feature} mạnh --- trong thực tế gần như không bao giờ đúng tuyệt đối.
  \item Vấn đề \textbf{zero frequency}: nếu một giá trị phân loại chưa xuất hiện trong tập huấn luyện, mô hình gán xác suất 0 và không dự đoán được (cần làm trơn Laplace trong thực hành).
\end{itemize}
\end{luuy}

\subsection{Ứng dụng phân loại Naive Bayes}

\begin{ynghia}[title={Dự đoán thời gian thực}]
Triển khai và tính toán nhanh $\Rightarrow$ phù hợp dự đoán real-time.
\end{ynghia}

\begin{ynghia}[title={Dự đoán đa lớp}]
Ước lượng xác suất hậu nghiệm của nhiều lớp biến mục tiêu.
\end{ynghia}

\begin{ynghia}[title={Phân loại văn bản}]
Phù hợp spam-filtering, phân tích cảm xúc (sentiment analysis).
\end{ynghia}

\begin{ynghia}[title={Hệ gợi ý (Recommendation system)}]
Kết hợp collaborative filtering: lọc thông tin chưa thấy, dự đoán người dùng có thích tài nguyên hay không.
\end{ynghia}
